\subsection{Field orientation relative to hippocampal anatomy}\label{supp:fem-orientation}
The fields in IXI002 follow neither the hippocampal long axis nor the local
laminar direction. Table~\ref{tab:fem-orientation} compares the carrier fields
$\mathbf E_1,\mathbf E_2$ and each montage's balanced-envelope-maximizing
direction $\mathbf n^*$ at $2+2$~mA.
The principal coordinate axes of the 1,568-node FreeSurfer hippocampal region
of interest (ROI) are PC1, the anterior-positive long axis; PC2, lateral-positive; and
PC3$={\rm PC1}\times{\rm PC2}$, normal to its best-fit plane. Angles are
unsigned, from $0$ to $90^\circ$, and nodes receive equal weight.
Weighting by field magnitude, or envelope maximum for $\mathbf n^*$,
changes these angular medians by at most about $1.5^\circ$.

\begin{table}[htbp]
\centering\small
\setlength{\tabcolsep}{3pt}
\begin{tabular}{@{}llrrrr@{}}
\toprule
Montage & Direction & PC1 & PC2 & PC3 & Laminar\\
\midrule
Violante & $\mathbf E_1$ & 71 (63--78) & 25 (17--33) & 79 (73--84) & 69 (53--80)\\
 & $\mathbf E_2$ & 58 (51--67) & 35 (27--42) & 83 (77--87) & 72 (57--81)\\
 & $\mathbf n^*$ & 66 (62--71) & 28 (22--34) & 80 (73--85) & 69 (53--80)\\
Yatsuda & $\mathbf E_1$ & 71 (64--78) & 25 (20--35) & 76 (69--84) & 74 (55--83)\\
 & $\mathbf E_2$ & 77 (71--83) & 21 (15--27) & 79 (70--85) & 74 (54--83)\\
 & $\mathbf n^*$ & 75 (66--83) & 23 (16--32) & 78 (70--85) & 74 (56--83)\\
\bottomrule
\end{tabular}
\caption{Angles in degrees, median (interquartile range). Coordinate-axis
comparisons use all 1,568 ROI nodes; laminar comparisons use the 934 nodes
where that direction is defined. Yatsuda denotes the FT7--F8, TP7--P8
rendering.}
\label{tab:fem-orientation}
\end{table}

The laminar direction is the normalized gradient of HippUnfold's inner--outer
Laplace coordinate \cite{dekraker2022}, sampled at the nearest voxel without
smoothing. HippUnfold 1.5.2-pre.2 processed the same T1 image; its container,
arguments and segmentation-model hash are recorded in
\path{docs/fem_ixi002_violante.md}. The coordinate has a $0.2$~mm grid
in T1 world coordinates. It defines directions at 934 nodes (59.6\% of the
ROI); 565 are interior, with their voxel and all six face neighbors inside
the Laplace domain. The laminar direction lies $33^\circ$ (21--54) from PC3
and $76^\circ$ (55--83) from PC2. Its median spread about its own principal
axis is $31^\circ$ (23--51), so one uniform direction cannot follow it
throughout the region.

A laminar-aligned detector samples smaller polarizations than an optimally
oriented one. In the Violante montage, median absolute projections of the
unit carrier directions onto the laminar direction are 0.36 and 0.31.
The map fractions in Table~\ref{tab:fem-laminar-fractions} are computed
node by node, each divided by that map's own maximum over orientation.
The quadratic and mean-square denominators therefore need not occur at
$\mathbf n^*$. These fractions describe geometry under a laminar-detector
hypothesis, before state-dependent detection or recurrent gain.

\begin{table}[htbp]
\centering\small
\begin{tabular}{@{}llrrr@{}}
\toprule
Montage & Nodes & Envelope & Quadratic beat & Mean square\\
\midrule
Violante & All defined (934) & 0.29 (0.11--0.57) & 0.11 (0.02--0.34) & 0.12 (0.04--0.33)\\
 & Interior (565) & 0.25 (0.10--0.54) & 0.08 (0.02--0.29) & 0.09 (0.03--0.31)\\
Yatsuda & All defined (934) & 0.25 (0.09--0.55) & 0.07 (0.01--0.32) & 0.08 (0.02--0.33)\\
 & Interior (565) & 0.25 (0.10--0.47) & 0.07 (0.02--0.25) & 0.08 (0.02--0.26)\\
\bottomrule
\end{tabular}
\caption{Laminar-projected map divided by its orientation maximum, median
(interquartile range). Restricting the Violante comparison to interior nodes
gives carrier and $\mathbf n^*$ angles to the laminar direction of
$71^\circ$, $74^\circ$, and $71^\circ$, respectively.}
\label{tab:fem-laminar-fractions}
\end{table}

The carrier directions themselves differ by a median $14.2^\circ$ (10.2--18.6) in the Violante montage and $7.5^\circ$ (4.2--11.4) in the Yatsuda rendering; their 95th percentiles are $29.2^\circ$ and $16.8^\circ$. The renderings also differ in electrode layout. With the tested 20-mm discs, the Violante e2 center is 9.4 and 18.4~mm from AF8 and Fp2, and e4 is 10.4 and 19.0~mm from P8 and PO8. These separations meet the 20-mm overlap criterion if those cap positions are occupied, for this layout and disc size.
